\begin{abstract}
On-policy self-distillation (OPSD) can inject a grounded hint into a teacher context and
pull the unhinted student distribution toward the resulting correction. This is attractive
because it converts semantic feedback into dense token-level supervision. However, the
student is asked to reproduce a distribution caused by information it does not observe.
When that information is recoverable from the deployed context, the update can compile a
useful behavior. When it is not recoverable, the same update can teach an unsupported
guess: the model learns what to say after an invisible intervention without learning how
to produce or justify that intervention. This proposal studies \emph{hinted supervised
fine-tuning} (hinted SFT), which directly trains a model to generate a grounded corrective
hint before acting and places the generated hint in the live inference context. The
proposal asks when this explicit intermediate improves causal faithfulness, long-horizon
credit localization, and retention relative to OPSD, direct SFT, and reward-only training.
\end{abstract}

\section{Research question and position}

\paragraph{Research question.}
For agent corrections that can be inferred from deployment-visible history, tool results,
or environment observations, does training the model to generate and use an explicit hint
produce more faithful and durable behavior than directly distilling a privileged
hint-conditioned distribution into an unhinted policy?

OPSD uses the same or a closely related policy as both student and teacher. The student
samples a state $x$ from its own rollout. The teacher additionally receives privileged
context $h$---for example, a verified demonstration or a short reminder at a localized
mistake---and produces a corrected token distribution. A KL loss moves the unhinted
student distribution toward that teacher distribution
\citep{shenfeld2026selfdistillation,zhao2026selfdistilledreasoner,cursor2026composer25}.
Composer 2.5 applies this idea to targeted textual feedback: a hint is inserted at a
problematic turn and the original-context policy is trained toward the hint-conditioned
policy \citep{cursor2026composer25}.

The proposal does not claim that OPSD is intrinsically non-causal or invalid. It isolates
a narrower condition:
\begin{center}
  \textbf{The desired correction must be supported by information available to the
  deployed student.}
\end{center}
If $h$ only restates evidence already in $x$, distillation may teach the model to attend
to that evidence without requiring it to print the hint. If $h$ contains a variable
unavailable from $x$, no policy conditioned only on $x$ can reliably reproduce different
teacher actions for different values of that hidden variable. Training pressure can then
produce an average, a dataset-prior guess, or a confident hallucination.

Hinted SFT makes the intermediate explicit. An offline reviewer first constructs a
validated hint $h^\star$ using information that will also be available at deployment. The
model is trained to emit $h^\star$ from $x$ before the corrected action. At inference, the
model generates $\hat h$, places it in its own context, and conditions later tokens on it.
The action distribution therefore changes through an observable intermediate generated
from observable evidence. The method is intended for diagnosis, planning, constraint
recovery, and local error correction; it is not a substitute for retrieving genuinely
missing information.
